% ================== Document class ==================
% Two-column layout, like a conference paper
\documentclass[10pt,twocolumn]{article}

% ================== Fonts and encoding ==================
\usepackage[T1]{fontenc}
\usepackage[utf8]{inputenc}
\usepackage{mathptmx}              % Times font (text and math), as in most papers
\usepackage{textcomp}
\usepackage{microtype}

% ================== Math ==================
\usepackage{amsmath,amssymb}

% ================== Page layout ==================
\usepackage[letterpaper,margin=0.75in,columnsep=0.25in]{geometry}
\setlength{\emergencystretch}{3em} % avoid overfull lines

% ================== Tables ==================
\usepackage{booktabs,array,tabularx}
\newcolumntype{L}{>{\raggedright\arraybackslash}X}
\renewcommand{\arraystretch}{1.15}

% ================== Sections ==================
% Section numbers are written by hand (e.g. "4.1 ..."),
% so automatic numbering is turned off.
\setcounter{secnumdepth}{-\maxdimen}

% ================== Links and colors ==================
\usepackage[dvipsnames]{xcolor}
\usepackage{xurl}
\usepackage[unicode]{hyperref}
\urlstyle{same}
\hypersetup{
  pdftitle={Beyond Accuracy: How Weight-Only Quantization Changes the
            Errors of Large Language Models in Reasoning and Factual Recall},
  pdfauthor={Abbass Aoun},
  colorlinks=true,
  linkcolor=Maroon,
  citecolor=Blue,
  urlcolor=Blue
}

% ================== Title ==================
\title{\Large\bfseries Beyond Accuracy: How Weight-Only Quantization Changes
the Errors of Large Language Models in Reasoning and Factual Recall}

\author{Abbass Aoun\\
\small \href{mailto:abbass.aoun@lau.edu}{abbass.aoun@lau.edu}\\[4pt]
\small COE 549: Advanced Large Language Models\\
\small Instructor: Dr.~Samer Saab Jr.}

\date{\small Submission date: November 3, 2026}

\begin{document}

% Title and abstract span both columns
\twocolumn[
\begin{@twocolumnfalse}
\maketitle
\begin{abstract}
\noindent Quantization is now the standard way to make large language
models (LLMs) small and cheap enough to run. There are many ways to do
it: plain round-to-nearest (RTN), methods that correct rounding errors
such as GPTQ, methods that protect the most important weights such as
AWQ, and methods that deal with outliers in the activations. Almost all
of them are judged the same way, by perplexity or by benchmark accuracy.
Recent work shows that this hides a lot. A quantized model can keep the
same accuracy while giving different answers to many questions, and some
abilities suffer much more than others. Multi-step reasoning, rare facts,
and answers the model was unsure about tend to break first, and they do
not all break in the same way. In this survey, we organize the
literature in two ways: by how each method handles rounding error, and
by what each study actually measures. We then review what is known about
reasoning, factual recall, changes in model behavior, failure modes, and
the effect of calibration data. We find three gaps. Most studies measure
\emph{how much} quality is lost, but few look at \emph{what kind} of
mistakes appear. RTN is rarely kept as a baseline next to GPTQ. And
reasoning and factual recall are almost never tested together in one
controlled setup. These gaps motivate our project: a controlled study of
the errors a small instruction-tuned model makes after being quantized
with RTN and with GPTQ.
\end{abstract}
\vspace{1.5em}
\end{@twocolumnfalse}
]

% ======================================================================
\section{1. Introduction}

Modern LLMs are expensive to run mainly because they are large. A
175-billion-parameter model stored in 16-bit floating point needs about
326~GB for its weights alone, far more than the 80~GB of a single
high-end GPU, so even inference requires several accelerators [2].
Quantization attacks this cost directly by storing each weight with
fewer bits. Moving from 16 to 4 bits cuts weight memory by roughly four
times, and because token-by-token generation is limited by how fast
weights can be read from memory rather than by arithmetic, it can also
make generation faster: GPTQ reports 3.25$\times$ to 4.5$\times$
speedups for a 3-bit OPT-175B [2]. This is why 4-bit weights are an
attractive choice for running open models on limited hardware: they give
the best zero-shot accuracy per bit [6] and are among the most efficient
formats for latency-sensitive serving [14].

The question that matters for practitioners is not whether a quantized
model is smaller, but whether it is still the same model. For most of
the short history of LLM quantization, the answer has been read off two
numbers: perplexity on a held-out corpus such as WikiText2, and average
accuracy on a handful of zero-shot multiple-choice benchmarks [2], [6],
[11]. When both stay close to the full-precision baseline, the quantized
model is treated as a drop-in replacement. Several recent studies
challenge that assumption. Dutta et al.~[26] show that compressed models
often stay within one or two points of baseline accuracy while flipping
up to 13.6\% of individual answers from correct to incorrect and back.
Jaiswal et al.~[13] show that compression methods reported as nearly
lossless in perplexity can still lose substantial accuracy on
knowledge-intensive question answering. Studies of mathematical
reasoning report losses that are much larger than on general
benchmarks, especially for small models and hard problems [11], [15],
[16], [18]. And interpretability work suggests that quantization fails
through qualitatively different mechanisms at different bit widths [23].

This survey takes that tension as its starting point. Its scope is
post-training quantization (PTQ) of decoder-only LLMs, with an emphasis
on weight-only quantization at 4 bits, which offers the best accuracy
per bit [6] and efficient deployment [14], and at 3 bits, where quality
starts to break down and most of the interesting failures appear [2],
[6]. We cover the main method families, but our real focus is
evaluation: what the literature has measured about the behavior of
quantized models, which abilities degrade first, how errors have been
characterized, and how the choice of calibration data shapes the result.
We draw on foundational methods papers, large benchmarking studies, and
a set of recent (2025 and 2026) papers that look beyond accuracy at
reasoning traces, failure modes, and factual knowledge.

The rest of the survey is organized as follows. Section~2 introduces the
notation and concepts needed to read the literature. Section~3 proposes
a taxonomy along two axes: how methods treat rounding error, and what
studies evaluate. Section~4 reviews the literature thematically,
comparing methods and findings rather than summarizing papers one at a
time. Section~5 examines datasets, metrics, and evaluation practice.
Section~6 discusses the main technical trade-offs. Section~7 identifies
open gaps, and Section~8 explains how those gaps shape our own project.
Section~9 concludes.

% ======================================================================
\section{2. Problem Formulation and Background}

\subsection{2.1 Uniform quantization and the RTN baseline}

A linear layer computes $Y = WX$, where $W$ is a weight matrix with
$d_{\mathrm{row}}$ rows (output channels) and $d_{\mathrm{col}}$ columns
(input channels), and $X$ holds the layer inputs. Uniform $b$-bit
quantization replaces each weight with one of $2^b$ evenly spaced
levels. For an asymmetric min-max grid, a group of weights with minimum
$w_{\min}$ and maximum $w_{\max}$ gets a scale and a zero-point:
\[
s = \frac{w_{\max} - w_{\min}}{2^b - 1}, \qquad
z = \mathrm{round}\!\left(-\frac{w_{\min}}{s}\right),
\]
\[
q = \mathrm{clip}\!\left(\mathrm{round}\!\left(\frac{w}{s}\right) + z,\
0,\ 2^b - 1\right), \qquad \hat{w} = s\,(q - z).
\]

How finely the scales are assigned matters as much as the number of
bits. With a per-tensor grid, the whole matrix shares one scale. With a
per-row (per-channel) grid, each output row gets its own scale. With a
group-wise grid, every block of $g$ weights in a row (usually $g = 128$)
gets its own scale and zero-point. Smaller groups give a closer fit but
cost a little extra storage, about 0.15 extra bits per weight for
$g = 128$ in GPTQ's setting [2].

Round-to-nearest (RTN) is the simplest method: it moves each weight to
the closest level on the grid, on its own, without looking at any data
or at the other weights. RTN is how LLM.int8() quantizes the weights
[1], and it is the standard baseline for large models in GPTQ [2] and in
S.~Li et al.~[11].

We use the common W/A/KV notation. W4A16 means 4-bit weights and 16-bit
activations (only the weights are quantized), W8A8 means 8-bit weights
and 8-bit activations, and KV4 means a 4-bit key-value cache. This
survey focuses mostly on W4A16 and W3A16.

\subsection{2.2 The layer-wise reconstruction objective}

Most accurate PTQ methods quantize a model one layer at a time [2], [3].
For each layer, they choose the quantized weights so that the layer's
\emph{output} stays close to the original, rather than the weights
themselves. This is measured on a small set of calibration inputs $X$:
\[
\hat{W} = \arg\min_{\hat{W}} \lVert WX - \hat{W}X \rVert_2^2 .
\]
For each row of $W$, this is a quadratic problem whose Hessian is
$H = 2XX^{\top}$. $H$ depends only on the inputs. Its diagonal shows how
large each input feature is, and its off-diagonal entries show which
input features move together. The Optimal Brain Quantizer (OBQ) [3],
which adapts the classic Optimal Brain Surgeon pruning method, uses
$H^{-1}$ to quantize one weight at a time. After each step, it adjusts
the weights that are not yet quantized so that they cancel the error
just introduced. GPTQ [2] keeps this idea but makes it fast enough for
models with hundreds of billions of parameters (Section~4.1).

The key point is that RTN and GPTQ can use exactly the same grid [2].
The only difference between them is whether rounding errors are
corrected.

\subsection{2.3 Outliers}

Quantization error depends on how wide a range the grid has to cover.
Dettmers et al.~[1] showed that once models reach about 6.7B parameters,
a few hidden dimensions start carrying very large activation values in
most layers and tokens. These outlier features stretch the quantization
range, and 8-bit activation quantization breaks down unless they are
handled separately.

S.~Li et al.~[11] measure the same effect with kurtosis, which tells how
many extreme values a distribution has. Weight kurtosis is moderate and
goes down as models grow (from 13.2 in OPT-1.3B to 5.2 in OPT-66B).
Activation kurtosis is much higher, in the hundreds or thousands, and it
goes up with model size. At the input of the FFN down projection in
LLaMA2-7B, it reaches about 154,000. This difference explains a pattern
seen across the literature: as models grow, weight-only quantization
gets easier, but activation quantization gets harder [11].

\subsection{2.4 Calibration data}

Methods that use data need a small calibration set to estimate $X$.
GPTQ uses it to build $H$ [2], while AWQ [4] and SmoothQuant [5] use it
to measure activation sizes. The usual setup, used by GPTQ and widely
reused since, is 128 random pieces of 2,048 tokens each from the C4 web
corpus [2]. This makes calibration a real design choice, not a small
detail: the inputs GPTQ sees decide which weights it thinks can make up
for the rounding errors of others. Section~4.5 looks at how much this
choice matters.

\subsection{2.5 Capability, distance, and flips}

Following Dutta et al.~[26], we separate two kinds of metrics.
Capability metrics, such as accuracy and perplexity, ask whether a model
is good. Distance metrics ask whether a compressed model still behaves
like the original. Examples are the KL divergence between the two
models' output distributions, and the share of answers that switch
between correct and incorrect. We call such a switch a \emph{flip}. In
this survey we separate \emph{harmful flips} (correct to incorrect) from
\emph{repair flips} (incorrect to correct). The two can balance out in
the accuracy score, even when the model's behavior has changed a lot
[26].

% ======================================================================
\section{3. A Taxonomy of the Literature}

Quantization papers are usually grouped by method. For our questions,
this hides an important split: papers that propose methods and papers
that evaluate them use different terms and ask different questions. We
therefore group the literature along two axes. The first groups methods
by how they handle rounding error (Table~\ref{tab:axisA}). The second
groups studies by what they actually measure (Table~\ref{tab:axisB}),
and this second axis turns out to be the more useful one.

\begin{table*}[t]
\centering
\small
\caption{Axis A: method families grouped by how they handle rounding
error.}
\label{tab:axisA}
\begin{tabularx}{\textwidth}{@{}p{0.15\textwidth} L p{0.27\textwidth} p{0.13\textwidth}@{}}
\toprule
\textbf{Family} & \textbf{Main idea} & \textbf{Methods} &
\textbf{Calibration data} \\
\midrule
Round-to-nearest & Each weight is rounded to the nearest grid level,
independently of the others & RTN; LLM.int8() weights [1]; [6] & Not
needed \\
Error compensation & After rounding some weights, the remaining ones are
updated to reduce the layer's output error, usually with second-order
information & OBQ [3], GPTQ [2], QEP [7] & Needed \\
Salience-aware scaling & Weight channels with large activations are
scaled up before rounding so they lose less precision & AWQ [4] &
Needed \\
Difficulty migration & Activation outliers are moved into the weights
with a per-channel scale that keeps the output unchanged & SmoothQuant
[5] & Needed \\
Rotation & Weights and activations are multiplied by orthogonal matrices
that spread outliers across channels & QuaRot, SpinQuant (see [8], [9]);
OneRot [8]; SpinOut [9]; DuQuant for TTS [10] & Only for learned
rotations \\
Mixed precision & The few outlier dimensions are kept in 16-bit while
the rest is quantized & LLM.int8() decomposition [1] & Not needed \\
\bottomrule
\end{tabularx}
\end{table*}

\begin{table*}[t]
\centering
\small
\caption{Axis B: studies grouped by what they measure, from the
shallowest level (1) to the deepest (5).}
\label{tab:axisB}
\begin{tabularx}{\textwidth}{@{}p{0.2\textwidth} L p{0.38\textwidth}@{}}
\toprule
\textbf{Level} & \textbf{What is measured} & \textbf{Representative
work} \\
\midrule
1. Language-modelling fit & Perplexity on WikiText2, C4, PTB & GPTQ [2],
QEP [7], most methods papers \\
2. Aggregate capability & Accuracy on zero-shot or multiple-choice
benchmarks, sometimes generative ones & S.~Li et al.~[11], Lee et
al.~[12], Kurtić et al.~[14], Dettmers and Zettlemoyer [6], Zhou et
al.~[25] \\
3. Targeted capability & Accuracy on one ability under controlled
conditions (math, factual recall, emergent abilities, calibration
domains) & Li et al.~[15], [16], P.~Liu et al.~[17], R.~Liu et
al.~[18], Jaiswal et al.~[13], Williams and Aletras [28], Y.~Liu et
al.~[29] \\
4. Behavioral distance & How far the quantized model's answers drift
from the baseline, regardless of accuracy & Dutta et al.~[26], Rababah
et al.~[27], output similarity in [14] \\
5. Mechanism and failure mode & Why examples fail, which components
break, what the failures look like & Wang et al.~[22], Zhou et al.~[23],
Chang et al.~[24], Lotfi et al.~[19], Lian et al.~[20], Alimaskina et
al.~[21] \\
\bottomrule
\end{tabularx}
\end{table*}

\subsection{3.1 Axis A: how a method handles rounding error}

Table~\ref{tab:axisA} groups the main method families. The first two
rows matter most for this survey. RTN and GPTQ can use exactly the same
grid, so comparing them tests one design choice: whether rounding errors
are left to add up or are actively cancelled. The other families mostly
target activation quantization, or use a different idea of which
weights matter.

\subsection{3.2 Axis B: what a study measures}

Table~\ref{tab:axisB} groups studies by how deep their evaluation goes.
Two points follow from it. First, most methods papers stop at levels 1
and 2, so when a method is called ``near-lossless'', this usually means
near-lossless in perplexity or multiple-choice accuracy. Second, the
studies that reach levels 4 and 5 rarely compare RTN and GPTQ under the
same setup (Lian et al.~[20] is an exception), and they usually study
either reasoning or factual knowledge, almost never both. The rest of
this survey is organized around these two points.

% ======================================================================
\section{4. Thematic Review of Prior Work}

\subsection{4.1 From independent rounding to error compensation}

The simplest method, RTN, works well at 8 bits. LLM.int8() [1] showed
that 8-bit weights rounded to nearest, together with careful handling of
activation outliers, lose no accuracy up to 175B parameters. Below 8
bits, things change. With 4-bit RTN, the WikiText2 perplexity of
OPT-175B rises from 8.34 to 10.54, which is worse than the 10$\times$
smaller full-precision OPT-13B. At 3 bits, RTN collapses to perplexities
in the thousands [2]. The reason is simple: RTN treats every weight on
its own, so rounding errors that point in the same direction just add up
in the layer output.

Error-compensating methods tackle this directly. OBQ [3] solved the
layer-wise problem of Section~2.2 almost exactly: it quantizes one
weight at a time, always picking the easiest one next, and updates the
remaining weights after each step. It matched the best post-training
methods of its time at 4 and 3 bits on ResNets. But its cost grows as
$O(d_{\mathrm{row}} \cdot d_{\mathrm{col}}^{3})$, which already meant
about an hour for a 25M-parameter ResNet-50. GPTQ [2] made three changes
that are each simple but together decisive. First, it quantizes columns
in one fixed order shared by all rows, so $H^{-1}$ is updated once per
column instead of once per weight. Second, it groups updates into blocks
of 128 columns, so it is not slowed down by memory bandwidth. Third, it
computes the needed rows of $H^{-1}$ in advance with a Cholesky
decomposition and 1\% dampening, which removes the numerical problems
that otherwise appear in billion-parameter layers. As a result, OPT-175B
can be quantized in about four GPU hours, with 4-bit perplexity almost
unchanged (8.34 to 8.37) and 3-bit perplexity still usable (8.68). The
fixed order is the interesting choice. GPTQ shows that OBQ's careful
ordering gains little on large layers: the weights with large errors
that it saves for last find few weights left to absorb their error.

Two results complicate the simple story that compensation always wins.
Dettmers and Zettlemoyer [6] ran more than 35,000 experiments with
quantization that uses no data. They found that 4-bit RTN-style
quantization with small blocks already gives the best accuracy per bit
on zero-shot tasks, and that block size and data type matter more than
most other tricks. So RTN with groups of 128 is a serious baseline, not
a straw man. Arai and Ichikawa [7] went further. In their per-channel
experiments, 4-bit GPTQ failed badly on Llama-3-8B (WikiText2
perplexity 147.9, against 8.54 for RTN), although it helped on most
other models. They also argue that layer-wise methods, including GPTQ,
ignore the error that builds up across layers. Their QEP correction
improves nearly every method they test (RTN, GPTQ, AWQ and QuIP), most
clearly at 2 bits. GPTQ partly addresses this by feeding each block the
outputs of blocks that are already quantized [2], but QEP shows that the
remaining error between layers is still large enough to matter.

The other method families start from a different idea of what goes
wrong. AWQ [4] argues that only about 1\% of weight channels are
salient, that these are best found by activation size rather than
weight size, and that scaling them up before rounding protects them
without any second-order correction. It reports better results than RTN
and similar or better results than GPTQ at 3 and 4 bits, and claims to
overfit less to the calibration data. Independent comparisons disagree
on which is better. Lee et al.~[12] find AWQ generally ahead of GPTQ at
4 bits on instruction-tuned models, while Kurtić et al.~[14] find GPTQ
ahead on real-world tasks, using MSE-optimal clipping and better
calibration data. This disagreement is useful in itself: rankings
between 4-bit methods depend on implementation details and evaluation
choices, so claims that one method is better should be read with care.

Activation quantization has its own line of work. LLM.int8() [1] keeps
the outlier dimensions in FP16, which preserves accuracy but gives no
speedup on small models. SmoothQuant [5] instead divides each activation
channel by a smoothing factor and multiplies the matching weights by the
same factor, so both become easy to quantize in INT8. It matches FP16
accuracy up to 530B parameters, with up to 1.56$\times$ speedup.
Rotation-based methods push this idea to 4-bit weights, activations and
KV cache: they multiply tensors by orthogonal matrices that spread
outliers across channels. Park et al.~[8] show that random rotations
(QuaRot) give results that vary noticeably between random seeds, and
that learning the rotation from a single perturbed sample (OneRot)
matches SpinQuant, which uses 800 samples, while reducing that variance.
Their journal extension, SpinOut [9], injects outliers only into the
layers ranked most sensitive and improves W4A4KV4 perplexity further.
Li [10] applies DuQuant-style rotation to the language model inside a
text-to-speech system. Under W4A8 quantization, speaker similarity and
perceived quality stay the same, but the character error rate rises: a
model can sound fine while saying the wrong thing. These families matter
for the wider picture, but because they target activations, their
evaluations rarely isolate what weight rounding alone does to a model's
answers.

\subsection{4.2 What quantization does to reasoning}

Across studies, multi-step reasoning is the most fragile ability.
S.~Li et al.~[11], in an evaluation of eleven model families, find
multi-step reasoning, especially math, among the most sensitive emergent
abilities. Zhou et al.~[25] fit task-stratified scaling laws over 293
GPTQ configurations of Qwen3 and Llama-3. They call reasoning
``precision-critical'': it is the ability most sensitive to bit width,
by about 40\% compared with the others, which they explain by errors
adding up across reasoning steps.

How large the effect is depends strongly on model size and task
difficulty. P.~Liu et al.~[17] find that 4-bit GPTQ keeps in-context
learning, chain-of-thought and instruction following on LLaMA models
from 7B to 65B, while 2-bit models collapse. Their GSM8K scores barely
change at 4 bits (13.1 to 13.5 for LLaMA-7B), but the full-precision
scores are so low that there is little to lose. This floor effect hides
small reasoning losses. Studies with stronger instruction-tuned models
see much more. Lee et al.~[12] report that GPTQ drops Llama-3.2-1B from
33.36 to 8.04 on GSM8K, while models of 70B and above barely change.
Li et al.~find that 4-bit AWQ and GPTQ cost Llama-3 models up to about
32\% relative accuracy on MATH, for the 1B model [15], and, in the
expanded study, that GPTQ lowers Qwen2.5-1.5B on GSM8K from 68.23 to
59.89 [16]. For long chain-of-thought reasoning models, R.~Liu et
al.~[18] find W8A8 almost lossless and 4-bit weights ``nearly
lossless'' (for example, a 2.1\% drop for AWQ on the 1.5B and 7B
models). However, losses grow sharply at lower bit widths, are
concentrated in the smaller models, and are up to about four times
larger on AIME than on GSM8K. The common thread is that quantization
hurts most for small models and hard, long problems, which is exactly
where the average benchmarks used by methods papers say the least.

Which kinds of errors appear is less clear. Li et al.~[15] label each
wrong answer by its first wrong step, using a judge model, a vote among
five LLMs and a human check of 2\% of cases. They sort errors into
conceptual, method, execution and reasoning errors. Conceptual errors
remain the largest group after quantization, and execution errors come
second. Multiplication is especially affected. Reasoning errors look
rare, but the authors explain this as a survivor effect: only the first
error is counted, and quantized models fail early. S.~Li et al.~[11], in
a small analysis of LLaMA2-70B, instead find wrong logic behind about
half of the math errors. The two analyses differ in model size, bit
width and categories, so they cannot be compared directly. Together,
they show that the answer depends on the setting, and that a controlled
comparison is missing. P.~Liu et al.'s case study [17] adds a
qualitative contrast: at 4 bits the model writes a fluent solution with
one wrong step, while at 2 bits its text is no longer a solution at all.

The newest work on reasoning models looks at the reasoning trace
itself. Lotfi et al.~[19] find that the main failure under aggressive
quantization is overthinking: the model reaches the correct answer
midway but does not commit to it. On MATH-500 with 3-bit AWQ, such
errors rise from 19 to 139 and make up 52\% of failures. Notably, the
per-token difference between the quantized and full-precision outputs
follows the full-precision model's next-token entropy ($\rho = 0.92$):
quantization noise does its damage where the original model was already
unsure. Lian et al.~[20] show that quantized reasoning models often keep
their accuracy but write much longer reasoning. At INT4 on Qwen3-4B,
reasoning grows by 12\% with GPTQ but by 42.5\% with RTN, and at INT3 by
276\% with AWQ and 292\% with HQQ, enough to cancel the expected
speedup. Alimaskina et al.~[21], working at 2 bits, describe
path-finding failures, where the trace loops before any answer appears,
and commitment failures, where an answer appears but the model keeps
checking it. They find that the loop rate predicts accuracy loss
($r = -0.85$). These studies agree that quantization changes how a
model reasons, not only whether it ends with the right number. Lian et
al.~also give one of the few direct RTN-versus-GPTQ comparisons in this
literature.

\subsection{4.3 What quantization does to factual knowledge}

Factual recall behaves differently from reasoning, and studies disagree
on how much it suffers. Studies that report average scores tend to find
small losses. Wang et al.~[22] test Llama3-8B, Qwen2.5-7B and
Qwen2.5-14B on the LRE fact-probing dataset. On Qwen, 4-bit methods
reduce recall by about 2.5 to 3 points, while 8-bit is almost lossless,
and the authors conclude that quantization is still practical for
knowledge-heavy use. Dutta et al.~[26] see only 2--4\% flips on
TriviaQA, against 12--30\% on GSM8K, for large models compressed mainly
with bitsandbytes. Other studies see much larger effects. Jaiswal et
al.~[13] report that even 8-bit GPTQ loses about 8--10\% on exact-match
factoid QA (FreebaseQA) for Vicuna models, while the same models look
much healthier on multiple-choice MMLU, because picking one of four
options is much easier than generating the exact entity. Their open-book
experiments add a useful distinction: when the needed knowledge is given
in the prompt, compressed models still retrieve and summarize well. So
what compression mainly damages is the knowledge stored in the weights.

The more interesting results are about which facts are lost, and how.
Wang et al.~[22] find that relations where the full-precision model is
not already near-perfect lose the most, so partly learned knowledge is
the most fragile. Their neuron-attribution analysis places the loss in
the last layers for Qwen and in the middle-to-late layers for Llama. In
two-hop questions, the first hop (recalling the bridge entity) drops by
up to 30.08\%, while the second hop drops by only 4.25\% on average.
Chang et al.~[24] show that the same inputs are hard for GPTQ, AWQ,
NormalFloat and EfficientQAT alike (average correlation 0.82 across 50
method pairs). Examples with small residual-stream values in late layers
suffer most, because errors built up across layers are amplified there.
On PopQA, mid-popularity entities lose more than 10\%, while very famous
and very obscure ones are hurt less. Zhou et al.~[23] give the clearest
picture of what the failures look like. At 4 bits, the correct answer
usually drops in rank but stays in the top five: this is signal
degradation, where the computation still works but is noisy. At 2 bits,
its rank falls into the thousands and the model outputs common tokens
such as ``the'': this is computation collapse, where key components
stop working. Patching in activations from the full-precision model
fixes the first failure but not the second.

Two features of this literature limit what we can conclude. First, most
factual studies use probing formats (fill-in-the-blank prompts scored by
the next token) instead of open questions answered in free text. This
suits interpretability, but it hides what an error looks like in
practice. Second, the results are inconsistent across models. Wang et
al.~report that 4-bit GPTQ collapses Qwen2.5-14B from 73.1\% to 25.2\%
while leaving the 7B model mostly intact, and that quantization
sometimes improves recall [22]. Because they use pre-quantized public
checkpoints, differences in calibration data and settings cannot be
ruled out as causes. Zhou et al.~[25] add one more hint: in their
scaling laws, memorization is about twice as sensitive to the size of
the calibration set as application tasks are. This suggests that
factual recall may depend on calibration choices more than average
scores show.

\subsection{4.4 Beyond accuracy: flips, distances and failure modes}

A separate line of work argues that accuracy is the wrong thing to look
at. Dutta et al.~[26] make this case most directly. Across six
quantization schemes and seven multiple-choice tasks, accuracy stays
within about 2\%, while flips reach 13.6\%; only 8-bit GPTQ keeps flips
negligible. Flips increase with lower bit widths, harder tasks and
smaller models, and they correlate strongly with the KL divergence
between output distributions (Spearman 0.981 on MMLU). Their explanation
of why flips balance out is useful. The answers that change are those
with a small top margin (the gap between the best and second-best
option). Wrong answers tend to have smaller margins, so they change more
often (for 4-bit GPTQ on Llama2-7B, 15.3\% of wrong answers changed,
against 9.3\% of correct ones). Rababah et al.~[27] reach a similar
conclusion with a different measure, correctness agreement, which is
always below accuracy and falls as precision drops. Kurtić et al.~[14]
compare the free-text outputs of quantized 70B and 405B models with the
BF16 outputs. They find an average ROUGE-L of about 0.56 but a
BERTScore of 0.93: the wording often changes even when the meaning stays
close. Perplexity hides the same effect. Dutta et al.~show that
quantization acts like roughly symmetric noise on token
log-probabilities, which leaves perplexity almost unchanged while
generation quality gets worse.

Distance metrics tell us that behavior changed, but, as Dutta et
al.~note, a flip is ``only a warning''. To know what changed, we need
failure modes. Here the literature offers several taxonomies, each tied
to one task: first-error categories for math solutions [15], signal
degradation versus computation collapse for factual recall [23],
overthinking for reasoning models [19], path-finding failures,
commitment failures and loops at 2 bits [21], and token inflation [20].
These taxonomies are valuable, but they cannot be compared with each
other: each was built for one type of task, mostly with one quantization
method, and none is applied to reasoning and factual recall side by
side. The flips literature, meanwhile, counts changes but does not
classify them, apart from a qualitative look at around ten MT-Bench
answers in [26].

\subsection{4.5 The role of calibration data}

If GPTQ corrects errors based on the statistics of its calibration
inputs, then the choice of calibration data should matter. The evidence
suggests that it matters little for general benchmarks at 4 bits, but a
lot for specialized tasks at 3 bits. Williams and Aletras [28], in the
most systematic study, compress LLaMA, Vicuna and OPT models with GPTQ,
SpQR, SparseGPT and Wanda, using five calibration sources and ten random
sets from each. Pruning turns out to be sensitive, but 4-bit GPTQ
changes by only about 0.1 to 0.4 points in average accuracy between
sources, and using more than about 128 calibration examples brings
little extra gain. Y.~Liu et al.~[29] test whether calibration data from
the same distribution as the test set helps, and conclude that it is
``not necessarily optimal'': the best calibration source depends on the
dataset and the method, although same-distribution data usually helped
in their Chinese cross-dataset tests.

The picture changes for low bits and specialized domains. R.~Liu et
al.~[18] calibrate 3-bit GPTQ on a 1.5B reasoning model with either
WikiText2 or self-generated math reasoning data. Math calibration raises
average accuracy from 28.3 to 38.1, and GSM8K from 57.5 to 75.7. Lian et
al.~[20] report that math chain-of-thought calibration raised INT3
accuracy from 28\% to 65\% and cut reasoning length from about 30k to
12k tokens. Arai and Ichikawa [7] find that GPTQ calibrated on PTB did
worse than plain RTN on WikiText2, and Kurtić et al.~[14] report that
calibrating 4-bit GPTQ on random tokens hurts accuracy. AWQ, which uses
calibration data only to measure activation sizes, is claimed to be less
prone to this kind of overfitting [4].

What these studies have not asked is what domain-specific calibration
costs elsewhere. Every experiment that shows math calibration helping
math reasoning measures math reasoning only. None checks whether the
same model gets worse at recalling facts, and none checks whether
calibration changes the kind of errors that remain.

% ======================================================================
\section{5. Datasets, Benchmarks and Evaluation Protocols}

The benchmarks used in the reviewed work fall into a few families
(Table~\ref{tab:bench}), and the family a study chooses largely decides
what it can find.

\begin{table*}[t]
\centering
\small
\caption{Benchmark families used in the reviewed literature.}
\label{tab:bench}
\begin{tabularx}{\textwidth}{@{}p{0.17\textwidth} L p{0.22\textwidth} p{0.2\textwidth}@{}}
\toprule
\textbf{Family} & \textbf{Typical datasets} & \textbf{Main metric} &
\textbf{Used in} \\
\midrule
Language modelling & WikiText2, C4, PTB & Perplexity & [2], [7]--[9],
[27], [28] \\
Multiple-choice and zero-shot & MMLU, ARC, PIQA, HellaSwag, Winogrande,
LAMBADA, StoryCloze & Accuracy (option log-likelihood, sometimes
length-normalized) & [2], [6]--[8], [11], [12], [26]--[28] \\
Generative reasoning (mostly math) & GSM8K, MATH / MATH-500, AIME, GPQA,
BBH & Exact match on the final answer & [12], [14]--[21], [25], [26] \\
Factual recall & TriviaQA, FreebaseQA, PopQA, Natural Questions, LRE,
TwoHop-Fact, ParaRel; TruthfulQA (truthfulness) & Exact match or
next-token accuracy & [12]--[14], [22]--[26] \\
Open-ended generation & MT-Bench, Arena-Hard & LLM-judge scores,
similarity to the baseline output & [12], [14], [26] \\
\bottomrule
\end{tabularx}
\end{table*}

The usual baselines are the 16-bit model and, in methods papers, RTN on
the same grid [2]. A clear trend in recent evaluation studies is that
RTN disappears. The analyses of reasoning errors [15], factual recall
[22], failure modes [23] and overthinking [19] compare GPTQ, AWQ or
other methods with each other and with full precision, but not with RTN.
This matters because RTN is the only baseline that isolates the effect
of error compensation.

Several weaknesses come up again and again in this literature. First,
perplexity and multiple-choice accuracy miss the damage that matters for
generation [13], [26]: scoring only the first generated token cannot
reveal a wrong calculation or an answer that names a nearby but wrong
entity. Second, weak baselines create floor effects. When a
full-precision model solves only 13\% of GSM8K, as in [17], a 4-bit
model has little room to lose, and the study will conclude that
reasoning is robust. Third, samples are sometimes small. Lian et al.~[20]
use around 100 questions per benchmark, and the qualitative analysis in
[26] looks at about ten MT-Bench answers. This is enough to show a
trend, but not to compare fine-grained categories. Fourth, LLM judges
make their own mistakes. Lee et al.~[12] note that GPT-4 sometimes
graded wrong MT-Bench answers as correct, and Li et al.~[15] report a
judge treating ``5.5'' and ``11/2'' inconsistently. Without answer
normalization and a human check, some reported ``flips'' are not real
errors at all.

Finally, several studies change the model or the quantization pipeline
in ways that blur the comparison. Wang et al.~[22] use public
pre-quantized checkpoints, so calibration data and settings differ
between methods and models. Li et al.~[15] fine-tune every model for
output format before evaluation. This makes step-level analysis
possible, but it means the tested models are no longer the original
ones. Park et al.~[8] show that results can vary between random seeds,
yet most studies report a single run. And some papers contain internal
inconsistencies, such as full-precision scores that differ between
tables in [15], or appendix results that contradict the main text in
[27]. None of this cancels the findings, but it argues for controlled
pipelines, repeated runs and clear answer normalization.

% ======================================================================
\section{6. Technical Trade-offs and Design Choices}

\textbf{Bit width against model size.} For a fixed memory budget, a
larger model at lower precision often beats a smaller model at higher
precision. Dettmers and Zettlemoyer [6] find that 4 bits is almost
always the best choice for zero-shot accuracy per total model bit, and
that the trend reverses at 3 bits. S.~Li et al.~[11] show a 3-bit
LLaMA2-70B beating a 16-bit LLaMA2-13B on RACE with the same 27~GB of
memory. The catch is that this trade-off is measured on average
benchmarks. Reasoning studies keep finding that small models [12], [15],
[18] and hard problems [16], [18] suffer most, so a 4-bit setting that
looks safe on average can still hide large losses on exactly the tasks
users care about.

\textbf{Memory against latency.} Weight-only quantization always saves
memory, but speed is a separate question. GPTQ's 3.25$\times$ to
4.5$\times$ speedups for OPT-175B come from reading fewer bytes during
token-by-token generation, which is limited by memory, not from faster
arithmetic [2]. LLM.int8(), in contrast, saves memory but is often
slower than FP16 on small models [1]. Kurtić et al.~[14] show that the
best format depends on how the model is served: W4A16 is the most
cost-efficient when requests are handled one at a time, while W8A8 gives
the highest throughput with asynchronous continuous batching. Reasoning
models add a hidden cost. If a quantized model writes many more tokens
(42.5\% more with RTN at INT4 on Qwen3-4B [20]), or loops until it runs
out of budget [21], the per-token speedup can vanish or even reverse.
Alimaskina et al.~[21] report that plain 2-bit inference was slower
end-to-end than FP16 in most of their failure regimes.

\textbf{Quantization quality against quantization cost.} RTN takes
seconds and needs no data. OBQ is accurate but works in practice only up
to about 100M parameters [2]. GPTQ sits in between: it takes minutes for
1--3B models and about four GPU hours for 175B [2]. QEP improves quality
further, especially at low bit widths, with a correction that costs less
than the quantization itself [7]. Quantization-aware training was the
only method that improved both accuracy and reasoning length in [20],
but at a much higher cost. Cheap repairs after quantization are another
option. Zhou et al.~[23] fix much of the 4-bit signal degradation
without any training, but cannot fix 2-bit collapse, and Li et al.~[15]
recover most of the lost math accuracy with a few minutes of targeted
DPO fine-tuning.

\textbf{Weights against activations.} Weight-only quantization gets
easier as models grow, because extreme weight values become rarer, while
activation quantization gets harder, because activation outliers grow
[11]. W4A8 is roughly the limit at which weight-activation quantization
still performs decently, and W4A4 breaks most models [11]. Methods that
handle activations (LLM.int8(), SmoothQuant, rotation methods) gain
speed and memory at the cost of extra machinery, so their evaluations
tell us less about what weight rounding alone does.

\textbf{Fitting the calibration data against generalizing beyond it.}
Methods that use data can tailor quantization to the inputs they expect.
At 3 bits, this helps a lot when the calibration data matches the task
[18], [20]. But GPTQ can also overfit to a poorly matched calibration
set and end up worse than RTN [7]. AWQ is designed to be less sensitive
[4], and QEP includes a strength parameter to limit overfitting [7]. The
open question is whether gains on the calibration domain come at a cost
on other abilities.

\textbf{Grid granularity against storage.} Smaller groups follow local
weight ranges more closely, and they are one of the few changes that
reliably improve accuracy per bit [6]. Grouping also makes extreme
settings usable: GPTQ with groups of 128 reaches about 2.2 bits per
weight with less than 1.5 points of perplexity loss on the largest
models [2]. The cost is extra storage for the scales and zero-points,
about 0.15 bits per weight at $g = 128$ [2], plus slightly more complex
kernels.

% ======================================================================
\section{7. Limitations and Open Gaps}

\textbf{Error compensation has been tested for how much it helps, not
for what it changes.} GPTQ's advantage over RTN is shown in perplexity
and benchmark accuracy [2], and Lian et al.~[20] show a difference in
reasoning length. No study we found asks whether correcting rounding
errors changes the kind of mistakes a model makes, for example whether
RTN produces more broken output or more arithmetic slips than GPTQ at
the same bit width. The existing error-type analyses [15], [19], [21],
[23] use GPTQ, AWQ or other methods, without an RTN control.

\textbf{Reasoning and factual recall are studied separately.} The
reasoning literature [15]--[21] and the factual-knowledge literature
[13], [22]--[24] use different models, formats and metrics. Zhou et
al.~[25] cover both in one scaling study, but only as average accuracy
within mixed ability groups. Dutta et al.~[26] report flip rates for
GSM8K and TriviaQA without looking at the errors, and Zhou et al.~[23]
check reasoning only briefly in an appendix. Whether quantization
damages these two abilities through the same kind of failure
(near-misses, collapse, or something else) has not been tested on the
same model under the same conditions.

\textbf{Error taxonomies are task-specific and rarely validated.} Each
failure-mode study builds its own taxonomy for one task. The labels come
from LLM judges, simple rules or manual inspection, and their
reliability is reported very differently: a five-model vote with a 2\%
human check in [15], but only a qualitative look at about ten examples
in [26]. Factual errors have no established taxonomy at the level of the
generated answer. The closest work describes them through internal
signals, such as the rank of the correct answer [23] or neuron
attributions [22].

\textbf{Small instruction-tuned models are under-studied.} Much of the
evidence comes from models of 7B parameters and above, from older base
models with low reasoning accuracy [17], or from long chain-of-thought
reasoning models [18]--[21]. Yet the studies that do include small
models find them the most sensitive [12], [15], [16], [18], and these
are also the models most likely to be quantized for edge devices.

\textbf{The cost of domain calibration on other abilities is unknown.}
Math-specific calibration clearly helps math at 3 bits [18], [20]. No
study we found has measured what it does to factual recall in the same
model, or whether it changes the types of errors that remain rather than
just their number.

\textbf{The link between uncertainty and failure is only partly
explored.} Dutta et al.~[26] show that multiple-choice answers with a
small margin flip most. Wang et al.~[22] show that partly learned
relations degrade most, and Lotfi et al.~[19] that quantization noise
concentrates where the full-precision model is uncertain. Chang et
al.~[24] similarly find that the full-precision model can predict which
inputs will break. Together, these results suggest that the original
model's confidence predicts which answers will break. Whether it also
predicts how they break, for instance whether low-confidence facts turn
into near-miss answers while confident ones collapse, has not been
studied.

% ======================================================================
\section{8. Implications for the Final Project}

The survey points to a specific and feasible niche. Our project will use
GPTQ [2] as its target method and RTN on the same grid as the control,
so that any difference can be traced to error compensation alone. It
will study a small instruction-tuned model, Qwen2.5-1.5B-Instruct,
where the literature expects the clearest effects and where published
GSM8K scores give a reference point [16]. It will test two abilities
with the same pipeline: multi-step math reasoning (GSM8K) and
open-domain factual recall from memory (TriviaQA, without context), both
answered in free text. Before the main analysis, we will reproduce
GPTQ's core comparison of WikiText2 perplexity for FP16, RTN and GPTQ
[2] to confirm that the quantization pipeline works.

Following Dutta et al.~[26], the unit of analysis will be the flip,
separating harmful from repair flips. We go one step further by
classifying the harmful flips into error types. For math, we adapt the
taxonomy of Li et al.~[15] and add a category for degenerate output such
as loops, which other studies report at low bit widths [17], [21]. For
facts, we use near-miss, unrelated, abstention and degenerate answers,
motivated by the difference between signal degradation and computation
collapse in Zhou et al.~[23]. A single LLM judge will label all flips,
and we will check it against a hand-labelled sample using Cohen's kappa.

Two secondary questions follow from the gaps above. The first is whether
calibrating GPTQ on math text (from the GSM8K training split, to avoid
test leakage), which helps math reasoning at low bit widths [18], [20],
comes at a cost to factual recall or changes the mix of errors. The
second, motivated by [19], [22], [26], is whether the full-precision
model's confidence predicts not only which answers flip but also what
kind of error each one becomes. Lessons from the evaluation literature
shape the design: the same group size for both methods [6], sanity
checks against known GPTQ failures [7], answer normalization before
counting flips [15], and a 3-bit setting in case 4-bit produces too few
flips to analyze.

% ======================================================================
\section{9. Conclusion}

The quantization literature has solved much of the engineering problem:
4-bit weight-only models are cheap, fast and, by the measures most
papers use, close to their full-precision versions. This survey has
argued that those measures are the wrong place to stop. Studies that
look past perplexity and multiple-choice accuracy consistently find that
quantization changes behavior unevenly: reasoning and rare or partly
learned facts suffer first, small models suffer most, many answers flip
even when accuracy holds, and failures take recognizable forms that
change with bit width, from noisy near-misses at 4 bits to collapse at 2
bits. At the same time, the literature rarely compares methods under
controlled conditions, rarely keeps RTN as a baseline, studies reasoning
and knowledge in isolation, and has no shared way of describing the
errors that quantization introduces. Characterizing those errors, and
how they depend on error compensation, task type, calibration data and
model confidence, is the natural next step.

% ======================================================================
\section{Statement on the Use of AI Tools}

Following the course guidelines and the instructor's permission, AI
tools were used to help draft and edit this survey, to check citations
against the original papers, and to simplify the wording. The author
selected and read the papers, verified every cited claim and reference,
and takes full responsibility for the content.

% ======================================================================
\begin{thebibliography}{99}
\small
\setlength{\itemsep}{2pt}

\bibitem{b1} T. Dettmers, M. Lewis, Y. Belkada, and L. Zettlemoyer,
``LLM.int8(): 8-bit matrix multiplication for transformers at scale,''
in \emph{Proc. NeurIPS}, 2022.

\bibitem{b2} E. Frantar, S. Ashkboos, T. Hoefler, and D. Alistarh,
``GPTQ: Accurate post-training quantization for generative pre-trained
transformers,'' in \emph{Proc. ICLR}, 2023.

\bibitem{b3} E. Frantar, S. P. Singh, and D. Alistarh, ``Optimal brain
compression: A framework for accurate post-training quantization and
pruning,'' in \emph{Proc. NeurIPS}, 2022.

\bibitem{b4} J. Lin \emph{et al.}, ``AWQ: Activation-aware weight
quantization for LLM compression and acceleration,'' in \emph{Proc.
MLSys}, 2024.

\bibitem{b5} G. Xiao, J. Lin, M. Seznec, H. Wu, J. Demouth, and S. Han,
``SmoothQuant: Accurate and efficient post-training quantization for
large language models,'' in \emph{Proc. ICML}, 2023.

\bibitem{b6} T. Dettmers and L. Zettlemoyer, ``The case for 4-bit
precision: k-bit inference scaling laws,'' in \emph{Proc. ICML}, 2023.

\bibitem{b7} Y. Arai and Y. Ichikawa, ``Quantization error propagation:
Revisiting layer-wise post-training quantization,'' in \emph{Proc.
NeurIPS}, 2025.

\bibitem{b8} S. Park, C. Kim, and K.-S. Chung, ``One-shot training for
reduction of performance variance in LLM rotation quantization,'' in
\emph{Proc. Int. Conf. Artif. Intell. Inf. Commun. (ICAIIC)}, 2026,
pp.~173--176, doi: 10.1109/ICAIIC68212.2026.11454275.

\bibitem{b9} S. Park and K.-S. Chung, ``SpinOut: Enhanced rotation-based
quantization for LLM by outlier injection,'' \emph{IEEE Access},
vol.~14, pp.~24082--24095, 2026, doi: 10.1109/ACCESS.2026.3664084.

\bibitem{b10} X. Li, ``Outlier-aware low-bit quantization for LLM-based
TTS: DuQuant MLP-only W4A8 quantization on TTS model,'' in \emph{Proc.
6th Int. Conf. Adv. Electr., Electron. Comput. Technol. (EECT)}, 2026,
doi: 10.1109/EECT68756.2026.11517795.

\bibitem{b11} S. Li \emph{et al.}, ``Evaluating quantized large language
models,'' in \emph{Proc. ICML}, 2024.

\bibitem{b12} J. Lee, S. Park, J. Kwon, J. Oh, and Y. Kwon, ``Exploring
the trade-offs: Quantization methods, task difficulty, and model size in
large language models from edge to giant,'' in \emph{Proc. IJCAI},
2025, pp.~8113--8121.

\bibitem{b13} A. Jaiswal, Z. Gan, X. Du, B. Zhang, Z. Wang, and Y. Yang,
``Compressing LLMs: The truth is rarely pure and never simple,'' in
\emph{Proc. ICLR}, 2024.

\bibitem{b14} E. Kurtić, A. N. Marques, S. Pandit, M. Kurtz, and D.
Alistarh, ``\,`Give me BF16 or give me death'? Accuracy-performance
trade-offs in LLM quantization,'' in \emph{Proc. ACL}, 2025,
pp.~26872--26886.

\bibitem{b15} Z. Li \emph{et al.}, ``Quantization meets reasoning:
Exploring LLM low-bit quantization degradation for mathematical
reasoning,'' arXiv:2501.03035, 2025.

\bibitem{b16} Z. Li \emph{et al.}, ``Quantization meets reasoning:
Exploring and mitigating degradation of low-bit LLMs in mathematical
reasoning,'' arXiv:2505.11574v4, 2026.

\bibitem{b17} P. Liu \emph{et al.}, ``Do emergent abilities exist in
quantized large language models: An empirical study,'' in \emph{Proc.
LREC-COLING}, 2024, pp.~5174--5190.

\bibitem{b18} R. Liu \emph{et al.}, ``Quantization hurts reasoning? An
empirical study on quantized reasoning models,'' in \emph{Proc. COLM},
2025.

\bibitem{b19} S. Lotfi, P. Kirichenko, S. Li, and Z. Liu, ``Quantized
reasoning models think they need to think longer, but they do not,'' in
\emph{Proc. NeurIPS}, 2026, to appear.

\bibitem{b20} X. Lian \emph{et al.}, ``Quantization inflates reasoning:
Token inflation as a hidden cost of low-bit reasoning models,''
arXiv:2606.25519, 2026.

\bibitem{b21} E. Alimaskina, D. Rudas, D. Shveykin, G. Molodtsov, P.
Vasiliev, and A. Beznosikov, ``Extreme low-bit inference in reasoning
models: Failure modes and targeted recovery,'' arXiv:2606.02011, 2026.

\bibitem{b22} Q. Wang \emph{et al.}, ``Through a compressed lens:
Investigating the impact of quantization on factual knowledge recall,''
in \emph{Proc. 6th Workshop on Trustworthy NLP (TrustNLP)}, 2026,
pp.~21--39.

\bibitem{b23} C. Zhou \emph{et al.}, ``From signal degradation to
computation collapse: Uncovering the two failure modes of LLM
quantization,'' in \emph{Findings of ACL}, 2026.

\bibitem{b24} T.-Y. Chang, M. Zhang, J. Thomason, and R. Jia, ``Why do
some inputs break low-bit LLM quantization?'' in \emph{Proc. EMNLP},
2025, pp.~3410--3429.

\bibitem{b25} C. Zhou \emph{et al.}, ``Task-stratified knowledge scaling
laws for post-training quantized large language models,'' in
\emph{Findings of ACL}, 2026.

\bibitem{b26} A. Dutta, S. Krishnan, N. Kwatra, and R. Ramjee,
``Accuracy is not all you need,'' in \emph{Proc. NeurIPS}, 2024.

\bibitem{b27} B. Rababah \emph{et al.}, ``The illusion of equivalency:
Statistical characterization of quantization effects in LLMs,''
arXiv:2607.08734, 2026.

\bibitem{b28} M. Williams and N. Aletras, ``On the impact of calibration
data in post-training quantization and pruning,'' in \emph{Proc. ACL},
2024, pp.~10100--10118.

\bibitem{b29} Y. Liu \emph{et al.}, ``Evaluating the generalization
ability of quantized LLMs: Benchmark, analysis, and toolbox,''
arXiv:2406.12928, 2024.

\end{thebibliography}

\end{document}
